% !TEX program = lualatex
% =============================================================================
%  Guia de Arquitetura — Hanaro — modelo AX Academy / IFAM
%
%  COMPILAR COM:  lualatex guia-arquitetura.tex     (duas vezes: sumário e páginas)
%
%  Versão resumida do documento completo (guia-usuario.tex): só a Parte I
%  (Arquitetura) e a Parte II (Inteligência Artificial). Os capítulos são os
%  mesmos arquivos de caps/, compartilhados com o documento completo; um ajuste
%  de texto neles vale para os dois. A Parte III (Referência Técnica) não entra.
%
%  Todo o estilo mora em axacademy-guia.sty. Aqui ficam só os dados do documento
%  e a ordem dos capítulos.
% =============================================================================
\documentclass[11pt,a4paper,oneside,openany]{book}
\usepackage{axacademy-guia}

% Documento completo (inclui a Parte III)? Os capítulos 1 e 2 são os mesmos nos dois
% documentos; só o que depende desta chave muda.
\newif\ifguiacompleto
\guiacompletofalse

% ---------- DADOS DESTE GUIA ----------
\guiatitulo{Guia de\\Arquitetura}
\guiasubtitulo{Sistema e agentes de inteligência artificial}
\guiasistema{Hanaro}
\guiatagline{}
\guiaequipe{equipe do Polo INOVA}
\guiaano{2026}
\guiaversao{0.1}
\guiacidade{Manaus, Amazonas --- Brasil}
\guiacliente{LG Electronics do Brasil}
\guiareferencia{}

% ---------- CONTATO (folha de rosto) ----------
% Campo vazio some da folha de rosto, com o ícone. O suporte do Hanaro é feito
% pelo próprio sistema (Capítulo 6), por isso não há e-mail.
\guiaemail{}
\guiasite{www.ifam.edu.br}
\guiaendereco{INOVA --- Polo de Inovação do IFAM\\
    Av. Gov. Danilo de Matos Areosa, 1975 --- Distrito Industrial I\\
    Manaus/AM --- CEP 69075-351}

% ---------- AVISO DA CAPA ----------
% Quebras explícitas: a caixa comporta três linhas curtas.
\guiaaviso{VERSÃO\\EM CONSTRUÇÃO}
\guiaescopo{Este guia apresenta, em linguagem acessível, como o Hanaro é feito e como
    funcionam os agentes de inteligência artificial. Ele não é o manual de uso: o passo
    a passo das telas está no Guia Rápido de Uso, e os diagramas e as especificações de
    engenharia estão na Documentação Técnica.}

% ---------- MARCA DO SISTEMA ----------
% O Hanaro tem o logotipo inteiro em arquivo (art/marca/hanaro-assinatura.pdf),
% então o ícone fica vazio e o nome não é composto em texto.
\guiamarcaicone{}
\guiamarcatexto{hanaro-assinatura}

% ---------- IMAGEM DA CAPA ----------
% Vazio: capa sem imagem. Para pôr uma, nomeie um arquivo de img/.
\guiacapaimagem{}

% ---------- TEXTO DA CONTRACAPA ----------
% Sai logo abaixo das marcas da AX Academy e do IFAM.
\guiasobreaequipe{Guia de Arquitetura do Hanaro: o sistema e os agentes de inteligência artificial. Baseado no estado do projeto em novembro de 2026}

\begin{document}

% ---------- ABERTURA ----------
\guiacapa
\guiafolhaderosto
\guiaavisodeuso

% ---------- HISTÓRICO DE REVISÕES ----------
\begin{guiahistorico}
    \guiarevisao{0.1}{Novembro de 2026}{AX Academy}{Parte I, Arquitetura, e Parte II, Inteligência Artificial.}
\end{guiahistorico}

\begin{remark}[info]
Este guia é a versão resumida da arquitetura do Hanaro e das funções de inteligência
artificial (assistente e tradução). O passo a passo das telas está no Guia Rápido de
Uso, e os diagramas e as especificações de engenharia estão na Documentação Técnica. O
assistente e a tradução são ligados por ambiente, aos poucos: se uma função não
aparece no seu Hanaro, ela ainda não foi ligada.
\end{remark}

\clearpage

\guiasobreax

\guiapagina{Sobre o sistema}{%
    O Hanaro é o sistema web de controle das perdas produtivas (\emph{scrap}) da
    Linha Final. Ele recebe o relatório exportado do GERP, consolida a base de scrap,
    publica o IF Cost no Dashboard e no Modo TV e ajuda analistas e gestores a
    investigar os casos relevantes e a produzir relatórios.

    Este guia tem duas partes. A Parte~I, \emph{Arquitetura}, mostra as peças
    do sistema, o caminho do dado do GERP ao Dashboard, os perfis de acesso e onde o
    sistema roda. A Parte~II, \emph{Inteligência Artificial}, explica o assistente de
    pré-análise, a consulta às análises aprovadas e a tradução de relatórios para o
    inglês, sempre com a decisão final nas mãos de uma pessoa.
}

\guiasumario

% ---------- PARTE I ----------
% Um capítulo só por parte; dentro dele, só seções.
\part{Arquitetura}
% =============================================================================
%  Capítulo 1 — Como o Hanaro é Feito (Parte I, Arquitetura)
%  Capítulo único da parte, só em seções. Explica, em linguagem de usuário, as
%  peças do sistema, o caminho do dado do GERP ao Dashboard, os perfis, onde o
%  sistema roda e o que fica registrado. Os detalhes de TI vão em caixas
%  "Para TI".
%
%  Estilos dos diagramas (arq ...) ficam aqui e valem para o resto do guia.
% =============================================================================
\tikzset{
    arq no/.style={draw=guiatinta,line width=0.8pt,fill=white,rounded corners=4pt,
        align=center,text width=2.7cm,inner sep=3.5pt,minimum height=0.95cm,
        font=\guiahead\fontsize{8.4}{10}\selectfont,text=guiatinta},
    arq ia/.style={arq no,draw=guiavermelho,line width=1.1pt},
    arq dado/.style={arq no,fill=guiaareia},
    arq fora/.style={arq no,draw=guiacinza,dash pattern=on 3pt off 2pt,text=guiacinza},
    arq grupo/.style={draw=guiaborda,line width=0.7pt,fill=guiapapel,rounded corners=6pt},
    arq titulo/.style={font=\guiahead\bfseries\fontsize{8.6}{10}\selectfont,
        text=guiagrafite,anchor=north west,inner sep=0pt},
    arq seta/.style={-{Stealth[length=5pt,width=4pt]},draw=guiatinta,line width=0.75pt,
        rounded corners=3pt},
    arq rotulo/.style={font=\guiahead\fontsize{7.4}{8.6}\selectfont,text=guiacinza,
        inner sep=1.5pt,fill=white},
}

\chapter{Como o Hanaro é Feito}\label{cap:arq}

\section{Público e escopo}

Esta parte do guia explica o Hanaro \emph{por dentro}, sem exigir conhecimento técnico.
Quem usa o sistema no dia a dia entende melhor o que ele faz, por que uma execução
demora alguns segundos para aparecer no Dashboard e o que acontece com os dados que
passam por ele. Os pontos que interessam à equipe de TI estão nas caixas
\textbf{Para TI}.

\begin{itemize}
    \item \textbf{Para quem é:} analistas, gestores e administradores do Hanaro.
    \item \textbf{O que cobre:} as peças do sistema, o caminho do dado do GERP ao
          Dashboard, os perfis de acesso, onde o sistema roda e o que fica registrado.
    \item \textbf{O que não cobre:} instalação, atualização e cópias de segurança.
\end{itemize}

\begin{remark}[info]
\textbf{Assistente e tradução: ativação gradual.} A pré-análise do assistente, a
consulta e a tradução de relatórios para o inglês dependem de um servidor de modelos
de IA. Elas são ligadas por ambiente, aos poucos. Se uma delas não aparece no seu
Hanaro, ela ainda não foi ligada; o restante do sistema funciona normalmente. A
Parte~II descreve essas funções.
\end{remark}

\section{Visão geral}

O Hanaro monitora os casos de scrap da fábrica e ajuda analistas e gestores a
investigar os eventos relevantes e a produzir relatórios sobre eles. Os dados vêm do
GERP; o sistema calcula os indicadores, dispara alertas, guarda a análise de cada
ocorrência, organiza os planos de ação e publica relatórios versionados.

A \autoref{fig:contexto} mostra quem fala com quem. Cada pessoa usa o sistema pelo
navegador, na intranet. O arquivo do GERP chega pelo robô de automação ou pelo envio
manual. As funções de inteligência artificial ficam em um servidor separado, e o
Hanaro só vai a ele quando precisa.

\begin{figure}[H]
\centering
\begin{tikzpicture}[x=1cm,y=1cm]
    \useasboundingbox (0,0.2) rectangle (16,-9.6);

    % ---- servidor do Hanaro ----
    \draw[arq grupo] (3.7,-0.3) rectangle (10.9,-7.5);
    \node[arq titulo] at (3.95,-0.5) {Servidor do Hanaro};
    \node[arq no] (web) at (5.4,-2.0) {Web\\[-1pt]{\fontsize{7}{8}\selectfont telas e gráficos}};
    \node[arq no] (api) at (5.4,-4.0) {API\\[-1pt]{\fontsize{7}{8}\selectfont regras e acesso}};
    \node[arq no,text width=3.4cm] (wrk) at (9.1,-4.0) {Worker\\[-1pt]{\fontsize{7}{8}\selectfont trabalho em segundo plano}};
    \node[arq dado,text width=4.0cm] (dados) at (6.0,-6.3)
        {PostgreSQL (dados)\\[-1pt]{\fontsize{7}{8}\selectfont e Redis (sessões e fila)}};

    % ---- servidor de IA (agente em cima, modelos no meio, tradução embaixo) ----
    \draw[arq grupo,draw=guiavermelho] (11.7,-0.3) rectangle (15.9,-7.5);
    \node[arq titulo,text=guiavermelho] at (11.95,-0.5) {Servidor de IA};
    \node[arq ia] (ag) at (13.8,-2.1) {Agente\\[-1pt]{\fontsize{7}{8}\selectfont pré-análise e consulta}};
    \node[arq ia] (md) at (13.8,-4.1) {Servidor de modelos\\[-1pt]{\fontsize{7}{8}\selectfont modelo de linguagem}};
    \node[arq ia] (tr) at (13.8,-6.1) {Tradução\\[-1pt]{\fontsize{7}{8}\selectfont português para inglês}};

    % ---- fora ----
    \node[arq no,text width=2.3cm] (pes) at (1.2,-2.0) {Analista\\Gestor\\Administrador};
    \node[arq no,text width=2.3cm] (rob) at (1.2,-4.0) {GERP e robô de automação};
    \node[arq fora,text width=2.0cm] (bcb) at (8.0,-8.8) {Banco Central\\[-1pt]{\fontsize{7}{8}\selectfont cotação PTAX}};
    \node[arq fora,text width=2.3cm] (tms) at (11.0,-8.8) {Microsoft Teams\\[-1pt]{\fontsize{7}{8}\selectfont avisos de alerta}};

    % ---- fluxos (retos, sem cruzar caixas) ----
    \draw[arq seta] (pes.east) -- node[arq rotulo,above=1pt] {navegador} (web.west |- pes.east);
    \draw[arq seta] (web) -- (api);
    \draw[arq seta] (rob.east) -- node[arq rotulo,above=1pt] {arquivo} (api.west |- rob.east);
    \draw[arq seta] (api) -- (wrk);
    \draw[arq seta] (api.south) -- (api.south |- dados.north);
    \draw[arq seta] ([xshift=0.2cm]wrk.south west) -- ([xshift=0.2cm]wrk.south west |- dados.north);
    \draw[arq seta] (wrk.east) -- ++(0.5,0) |- (ag.west);
    \draw[arq seta] (wrk.east) -- ++(0.5,0) |- (tr.west);
    \draw[arq seta] (ag) -- (md);
    \draw[arq seta] (tr) -- (md);
    \draw[arq seta] ([xshift=-0.6cm]wrk.south) -- ([xshift=-0.6cm]wrk.south |- bcb.north);
    \draw[arq seta] ([xshift=1.6cm]wrk.south) -- ([xshift=1.6cm]wrk.south |- tms.north);
\end{tikzpicture}
\caption{As peças do Hanaro e como se ligam. Linhas tracejadas: serviços externos.
Os serviços de IA são opcionais.}
\label{fig:contexto}
\end{figure}

\section{As peças do sistema}

Cada peça tem uma tarefa. A \autoref{tab:pecas} diz o que ela faz e em que situação o
usuário percebe a sua presença.

\begin{table}[H]
\caption{Peças do Hanaro, o que fazem e como o usuário as percebe.}\label{tab:pecas}
\guiazebra
\small
\noindent\begin{tabularx}{\textwidth}{@{}L{2.9cm}X L{4.3cm}@{}}
    \toprule
    \guiacabtabela
    \thead{Peça} & \thead{O que faz} & \thead{Quando você nota} \\
    \midrule
    \textbf{Web} & Mostra as telas e os gráficos no navegador. & Tudo o que você vê e
        clica \\
    \textbf{API} & Confere o que foi pedido, aplica as regras de negócio e responde. Os
        pedidos demorados vão para a fila. & Mensagens de erro de formulário; envio do
        relatório \\
    \textbf{Worker} & Faz o trabalho pesado em segundo plano: processa o arquivo do GERP,
        calcula os alertas, busca a cotação do dólar, exporta documentos e fala com os
        serviços de IA. & Status \emph{Na fila} e \emph{Concluído} em Execuções \\
    \textbf{PostgreSQL} & Guarda os dados do sistema: ocorrências, análises,
        relatórios e usuários. & Nada muda quando o sistema reinicia \\
    \textbf{Redis} & Guarda as sessões, a fila de trabalho e o cache do Dashboard. &
        Sessão que expira; Dashboard rápido \\
    \textbf{Automação do GERP} & Lê o arquivo exportado do GERP, normaliza, classifica e
        envia à API. & Execuções com origem \emph{GERP} \\
    \textbf{Agente} & Prepara pré-análises e responde às perguntas da Consulta, a
        partir das análises já aprovadas. & Selo \emph{Pré-análise pronta}; menu
        Assistente \\
    \textbf{Tradução} & Traduz o texto de um relatório do português para o inglês. &
        Opção \emph{Traduzir o texto para inglês} \\
\end{tabularx}
\end{table}

\begin{remark}[info]
\textbf{Para TI.} O agente e a tradução são serviços separados, com contêiner e chave
de acesso próprios. O Hanaro chama os dois, e eles só respondem: o agente nunca chama o
Hanaro de volta.
\end{remark}

\section{Do GERP ao Dashboard}

O caminho de um arquivo do GERP até o indicador na tela tem seis etapas
(\autoref{fig:caminho}). Entender a ordem explica por que o Dashboard só muda depois
que a execução está \emph{Publicada}.

\begin{figure}[H]
\centering
\begin{tikzpicture}[x=1cm,y=1cm,
        arq no/.append style={text width=4.3cm,minimum height=1.05cm}]
    \useasboundingbox (0,0.1) rectangle (16,-10.8);

    \foreach \i/\y/\titulo/\texto in {
        1/-0.6/{GERP}/{O relatório \emph{Other Account Transaction Text} é exportado em texto separado por tabulação.},
        2/-2.5/{Entrada do arquivo}/{O robô de automação lê o arquivo, normaliza, classifica e envia à API. O envio manual, na tela Execuções, faz o mesmo caminho.},
        3/-4.4/{API}/{Confere a forma da carga, aceita o pedido e o coloca na fila. A resposta é imediata; o processamento vem depois.},
        4/-6.3/{Worker}/{Compara o arquivo com o que já existe, grava as transações, atualiza os números do Dashboard e prepara o snapshot, tudo de uma só vez.},
        5/-8.2/{Execução publicada}/{O status da execução passa a \emph{Concluído} e o snapshot a \emph{Publicado}.},
        6/-10.0/{Dashboard e Modo TV}/{Leem os números já calculados. Os alertas avaliam a base e avisam no Hanaro e no Teams.}}
    {
        \node[arq no,anchor=west] (n\i) at (0.9,\y) {\titulo};
        \node[guia marcador,minimum size=13pt,font=\guiahead\bfseries\fontsize{7.5}{7.5}\selectfont]
            at (n\i.west) {\i};
        \node[font=\guiahead\fontsize{8.4}{10.4}\selectfont,text width=9.6cm,anchor=west,
              text=guiatinta,align=left] at (5.9,\y) {\texto};
    }
    \foreach \a/\b in {1/2,2/3,3/4,4/5,5/6}
        \draw[arq seta] (n\a.south) -- (n\b.north);
\end{tikzpicture}
\caption{Do GERP ao Dashboard, em seis etapas.}
\label{fig:caminho}
\end{figure}

Três regras do sistema decorrem desse caminho e valem a pena conhecer:

\begin{itemize}
    \item \textbf{O dólar é congelado na chegada.} O valor em US\$ usa a cotação PTAX
          do dia em que o dado chega ao Hanaro, obtida do Banco Central. Ele só muda
          em um recálculo manual. Não existe valor de dólar digitado à mão.
    \item \textbf{O que chegou não é apagado.} A ocorrência é a identidade do negócio.
          Se o conteúdo muda no arquivo seguinte, o Hanaro registra outra versão da
          transação; o que deixa de aparecer na extração passa a constar como ausente,
          sem sumir do histórico.
    \item \textbf{O Dashboard lê números prontos.} Ele consulta indicadores já
          calculados, e não as transações uma a uma. Por isso responde rápido, e por
          isso só reflete uma execução depois que o snapshot é publicado.
\end{itemize}

\begin{remark}[info]
\textbf{Para TI.} A escrita acontece no worker e a leitura, em projeção: transações,
projeção do Dashboard e snapshot ativo entram no mesmo commit. A reconciliação é por
partição (organização e dia). \ifguiacompleto As sequências completas estão na Parte~III.\fi
\end{remark}

\section{Perfis, sessão e segurança}

O Hanaro tem três perfis de conta. A \autoref{tab:perfis} resume o que cada um
encontra no menu. O acesso é sempre por usuário e senha, cadastrados pelo
administrador.

\begin{table}[H]
\caption{Perfis de conta e o que cada um usa.}\label{tab:perfis}
\guiazebra
\small
\noindent\begin{tabularx}{\textwidth}{@{}L{3.0cm}X@{}}
    \toprule
    \guiacabtabela
    \thead{Perfil} & \thead{O que usa} \\
    \midrule
    \textbf{Analista} & Dashboard, Execuções, Base de Scrap, Relatórios, Alertas e
        Configurações. No Assistente: as pré-análises e a Consulta (quando ligada). \\
    \textbf{Gestor} & O mesmo do analista, e ainda a \emph{Memória} e o
        \emph{Desempenho} do Assistente. \\
    \textbf{Administrador} & Uso da plataforma, Execuções, Suporte, Alertas,
        Integrações e a \emph{Memória do assistente}. O administrador com permissão de
        superusuário também vê Usuários e os registros de atividade (Logs). Não trabalha nas
        pré-análises. \\
\end{tabularx}
\end{table}

\begin{itemize}
    \item \textbf{Sessão.} A sessão fica guardada no navegador, em um cookie que o
          código da página não consegue ler. Cada gravação exige ainda um código de
          confirmação do próprio sistema, o que impede que outro site faça pedidos em
          nome do usuário.
    \item \textbf{Robô do GERP.} Não usa sessão de pessoa: envia a carga com uma chave
          de API que só permite criar registros de scrap.
    \item \textbf{Serviços de IA.} O Hanaro os chama com chaves próprias, e eles só
          respondem ao Hanaro.
\end{itemize}

\begin{remark}
Quando a sessão expira, o Hanaro abre a janela de acesso sobre a tela atual. Entre
de novo e continue de onde parou; o que já foi gravado não se perde.
\end{remark}


\section{Onde roda}

O Hanaro roda na intranet da empresa, em contêineres Docker, e \textbf{não é publicado
na internet}. Só o servidor web (nginx) abre uma porta; o banco de dados e o Redis
ficam em uma rede interna, sem saída para fora. A tela e a API usam o mesmo endereço, e
por isso não há troca de domínio entre elas.

A \autoref{fig:implantacao} mostra a divisão em servidores e redes. O servidor de IA
costuma ser outra máquina, com recursos próprios para os modelos.

\begin{figure}[H]
\centering
\begin{tikzpicture}[x=1cm,y=1cm]
    \useasboundingbox (0,0.2) rectangle (16,-8.2);

    \node[arq no,text width=2.4cm] (nav) at (1.4,-2.9) {Navegadores};

    \draw[arq grupo] (3.2,-0.3) rectangle (10.9,-7.7);
    \node[arq titulo] at (3.45,-0.5) {Servidor do Hanaro (Docker Compose)};
    \draw[arq grupo,fill=white] (3.5,-1.1) rectangle (10.6,-4.1);
    \node[arq titulo,text=guiacinza] at (3.7,-1.25) {rede proxy};
    \node[arq no,text width=1.4cm] (ngx) at (4.5,-2.9) {nginx\\[-1pt]{\fontsize{7}{8}\selectfont porta 4203}};
    \node[arq no,text width=1.4cm] (bck) at (6.9,-2.9) {backend\\[-1pt]{\fontsize{7}{8}\selectfont API}};
    \node[arq no,text width=1.4cm] (wk) at (9.6,-2.9) {worker};
    \draw[arq grupo,fill=white] (3.5,-4.6) rectangle (10.6,-7.4);
    \node[arq titulo,text=guiacinza] at (3.7,-4.75) {rede data};
    \node[arq dado,text width=2.0cm] (pg) at (6.9,-6.4) {PostgreSQL};
    \node[arq dado,text width=1.4cm] (rd) at (9.6,-6.4) {Redis};

    \draw[arq grupo,draw=guiavermelho] (11.7,-0.3) rectangle (15.9,-7.7);
    \node[arq titulo,text=guiavermelho] at (11.95,-0.5) {Servidor de IA};
    \node[arq ia,text width=3.4cm] (ag) at (13.8,-2.2) {agent-service\\[-1pt]{\fontsize{7}{8}\selectfont porta 8200 e banco próprio}};
    \node[arq ia,text width=3.4cm] (md) at (13.8,-4.3) {Servidor de modelos\\[-1pt]{\fontsize{7}{8}\selectfont Ollama ou equivalente}};
    \node[arq ia,text width=3.4cm] (tr) at (13.8,-6.4) {translation-service\\[-1pt]{\fontsize{7}{8}\selectfont porta 8100}};

    \draw[arq seta] (nav.east) -- node[arq rotulo,above=1pt] {HTTP} (ngx.west);
    \draw[arq seta] (ngx) -- node[arq rotulo,above=1pt] {/api} (bck);
    \draw[arq seta] (bck.south) -- (pg.north);
    \draw[arq seta] (wk.south) -- (rd.north);
    \draw[arq seta] (wk.east) -- ++(0.9,0) |- (ag.west);
    \draw[arq seta] (wk.east) -- ++(0.9,0) |- (tr.west);
    \draw[arq seta] (ag) -- (md);
    \draw[arq seta] (tr) -- (md);
\end{tikzpicture}
\caption{Servidores, redes e portas. Os dois blocos de IA são opcionais e rodam, em
geral, em outra máquina.}
\label{fig:implantacao}
\end{figure}

\begin{remark}[info]
\textbf{Para TI.} O worker é o único contêiner do Hanaro que também sai para fora da
rede: para o Teams, o Banco Central e os serviços de IA. A porta 8200 do agente e a
8100 da tradução devem ficar fechadas, exceto para o servidor do Hanaro.
\end{remark}

\section{Registro e rastreabilidade}

O sistema guarda o histórico do que acontece, para que um número possa ser explicado
depois.

\begin{itemize}
    \item \textbf{Versões que não mudam.} Relatórios publicados, análises de caso,
          metas, produção e IF Cost são versionados. Um relatório publicado aponta
          para um snapshot selado com um código de verificação (\emph{hash}): o número
          que saiu no relatório continua o mesmo, mesmo que a base seja atualizada.
    \item \textbf{Tarefas internas com nova tentativa.} Alertas, cartões do Teams,
          exportações e o envio ao agente nascem de registros gravados junto com a
          mudança que os causou. Se o serviço de destino falha, o Hanaro tenta de novo,
          com espera crescente, e nada se perde no caminho.
    \item \textbf{Registro de atividade.} As gravações de acesso, scrap, alertas e
          relatórios ficam registradas. A exportação desses registros é restrita ao
          superusuário.
    \item \textbf{Nada da IA vira dado sem decisão humana.} Uma proposta do assistente
          só entra em uma análise, relatório ou tarefa depois que uma pessoa aceita
          (Parte~II).
\end{itemize}

\section{Se algo falhar}

A \autoref{tab:falhas} reúne situações comuns, o que significam e o que fazer. Para
qualquer outra, avise a equipe de TI.

\begin{table}[H]
\caption{Situações comuns, o que significam e o que fazer.}\label{tab:falhas}
\guiazebra
\small
\noindent\begin{tabularx}{\textwidth}{@{}L{3.5cm}X L{4.6cm}@{}}
    \toprule
    \guiacabtabela
    \thead{O que você vê} & \thead{O que significa} & \thead{O que fazer} \\
    \midrule
    A execução fica \emph{Na fila} & O worker ainda não pegou o arquivo, ou está ocupado
        com outro. & Aguarde alguns instantes e atualize; se não mudar, avise a TI \\
    O Dashboard não mudou depois do envio & O snapshot da execução ainda está \emph{Não
        publicado}; o Dashboard mostra a execução anterior. & Confira o status em
        Execuções \\
    A tela pede o acesso de novo & A sessão expirou. & Entre de novo; nada se perde \\
    O assistente aparece como desligado & O servidor de IA está desligado ou fora do ar
        neste ambiente. & Siga o trabalho sem ele; as pré-análises já criadas continuam
        disponíveis \\
    A tradução não termina & O serviço de tradução está indisponível ou sem worker para
        a fila. & Desligue a tradução e emita em português, ou tente de novo \\
    O aviso não chegou ao Teams & A entrega falhou e será tentada de novo, com espera
        crescente. & Aguarde; se persistir, avise a TI \\
\end{tabularx}
\end{table}

\fimdecapitulo


% ---------- PARTE II ----------
\part{Inteligência Artificial}
% =============================================================================
%  Capítulo 2 — Inteligência Artificial (Parte II)
%  Capítulo único da parte. Seções: o que a IA faz e não faz, como nasce uma
%  pré-análise (BPMN), seis passos com a tela (lista do Assistente, leitura,
%  decisão, Consulta, Memória e Desempenho, tradução), selos e rotas, a
%  tradução por dentro, boas práticas e mensagens.
%
%  As telas de img/ são ESQUEMAS provisórios (rodapé "CAPTURA PENDENTE"): trocar
%  pelas capturas reais, na mesma resolução, e recalibrar os marcadores x/y.
% =============================================================================
\chapter{Inteligência Artificial}\label{cap:ia}

\section{O que a IA faz e o que não faz}

O Hanaro usa inteligência artificial em duas frentes: o \textbf{assistente}, que
prepara pré-análises e responde a perguntas sobre casos já tratados, e a
\textbf{tradução} dos relatórios para o inglês. Em ambas, a IA \emph{propõe} e uma
pessoa \emph{decide}. A \autoref{tab:funcoes} resume as funções.

\begin{table}[H]
\caption{Funções de IA do Hanaro.}\label{tab:funcoes}
\guiazebra
\small
\noindent\begin{tabularx}{\textwidth}{@{}L{2.7cm}X L{3.4cm}@{}}
    \toprule
    \guiacabtabela
    \thead{Função} & \thead{O que faz} & \thead{Quem usa} \\
    \midrule
    \textbf{Pré-análise} & Lê os fatos de um caso, busca casos parecidos na memória e
        propõe os campos da análise e as tarefas. & Analista e Gestor \\
    \textbf{Consulta} & Responde a uma pergunta em texto livre usando só análises já
        aprovadas, sempre com as fontes. & Analista e Gestor, quando ligada \\
    \textbf{Memória} & Reúne as análises aprovadas que servem de exemplo ao
        assistente. & Gestor e Administrador \\
    \textbf{Desempenho} & Mostra o que as pessoas decidiram sobre as propostas. &
        Gestor e Administrador \\
    \textbf{Tradução} & Traduz o texto de um relatório do português para o inglês
        antes do download. & Quem emite o documento \\
\end{tabularx}
\end{table}

\begin{remark}
\textbf{A IA nunca decide sozinha.} Nada do que o assistente propõe vira análise,
relatório ou tarefa até que uma pessoa aceite. O assistente também não altera dados do
Hanaro por conta própria e não chama o Hanaro: é o Hanaro que lhe entrega o caso e
busca a resposta.
\end{remark}

O que o assistente \textbf{não} faz: não afirma uma causa quando não há precedente (só
traz os fatos, o checklist 4M e perguntas para quem analisa); não cita fonte, número ou
data que não estejam nos dados; e não repete como certeza o que é hipótese. Toda causa
proposta é marcada como \emph{hipótese}, para ser confirmada por quem conhece a linha.

\begin{remark}[info]
\textbf{Ativação gradual.} O assistente, a Consulta e a tradução dependem do servidor
de IA e são ligados por ambiente. Quando estão desligados, o Hanaro avisa na própria
tela, as pré-análises já criadas continuam disponíveis e o restante do sistema segue
normal.
\end{remark}

\section{Como nasce uma pré-análise}

A pré-análise começa em uma \textbf{regra de alerta} configurada para pedir ajuda do
assistente. Quando a regra dispara sobre uma ocorrência, o Hanaro entrega o caso ao
agente e aguarda. O agente reúne os fatos, procura casos parecidos na memória e prepara
uma proposta, que passa por uma conferência antes de chegar à pessoa. A
\autoref{fig:ia} mostra o fluxo; os números nos quadros vermelhos levam aos passos
deste capítulo.

\begin{figure}[H]
\centering
\begin{tikzpicture}[x=1cm,y=1cm,
        bpmn tarefa/.append style={text width=4.0cm,minimum height=0.8cm}]
    \useasboundingbox (0,0) rectangle (16,-13.7);
    \bpmnraia{0}{5}{13.7}{Analista ou Gestor}
    \bpmnraia{5}{5.6}{13.7}{Hanaro}
    \bpmnraia{10.6}{5.4}{13.7}{Agente}

    % ---- Hanaro ----
    \node[bpmn inicio] (ini) at (7.8,-1.25) {};
    \node[bpmn rotulo,anchor=west] at ([xshift=4pt]ini.east) {Alerta dispara};
    \node[bpmn tarefa] (h1) at (7.8,-2.4) {Regra de alerta pede pré-análise};
    \node[bpmn tarefa] (h2) at (7.8,-3.7) {Entregar o caso ao agente};
    \node[bpmn tarefa] (h3) at (7.8,-8.0) {Mostrar o selo \emph{Pré-análise pronta}};
    \node[bpmn tarefa] (h4) at (7.8,-11.7) {Criar análise, relatório e tarefas aceitos};

    % ---- agente ----
    \node[bpmn tarefa] (a1) at (13.3,-3.7) {Reunir os fatos e buscar casos parecidos};
    \bpmngateway{g1}{13.3,-5.0}
    \node[bpmn rotulo,anchor=south west,inner sep=0pt] at ([xshift=2pt,yshift=-1pt]g1.north east) {Há precedente?};
    \node[bpmn tarefa,text width=2.3cm,minimum height=1.5cm] (a2) at (11.95,-6.5)
        {Adaptar campos e tarefas dos casos anteriores};
    \node[bpmn tarefa,text width=2.3cm,minimum height=1.5cm] (a3) at (14.65,-6.5)
        {Só fatos, checklist 4M e perguntas};
    \node[bpmn tarefa,text width=4.3cm] (a4) at (13.3,-8.0)
        {Conferir: recusar fonte ou número inventado};
    \node[bpmn tarefa] (a5) at (13.3,-11.7) {Guardar a análise aprovada na memória};
    \node[bpmn fim] (fim) at (13.3,-13.0) {};

    % ---- pessoa ----
    \bpmnpasso{p1}{2.5,-9.4}{1}{Abrir a lista do Assistente}
    \bpmnpasso{p2}{2.5,-10.55}{2}{Ler a pré-análise}
    \bpmnpasso{p3}{2.5,-11.7}{3}{Aceitar, editar ou rejeitar}

    % ---- fluxos ----
    \draw[bpmn fluxo] (ini) -- (h1);
    \draw[bpmn fluxo] (h1) -- (h2);
    \draw[bpmn fluxo] (h2) -- (a1);
    \draw[bpmn fluxo] (a1) -- (g1);
    \draw[bpmn fluxo] (g1.west) -| node[bpmn rotulo,anchor=north east,pos=0.15] {Sim} (a2.north);
    \draw[bpmn fluxo] (g1.east) -| node[bpmn rotulo,anchor=north west,pos=0.15] {Não} (a3.north);
    \draw[bpmn fluxo] (a2.south) -- (a2.south |- a4.north);
    \draw[bpmn fluxo] (a3.south) -- (a3.south |- a4.north);
    \draw[bpmn fluxo] (a4.west) -- (h3.east);
    \draw[bpmn fluxo] (h3.south) |- (p1.east);
    \draw[bpmn fluxo] (p1) -- (p2);
    \draw[bpmn fluxo] (p2) -- (p3);
    \draw[bpmn fluxo] (p3.east) -- (h4.west);
    \draw[bpmn fluxo] (h4.east) -- (a5.west);
    \draw[bpmn fluxo] (a5) -- (fim);
\end{tikzpicture}
\caption{Da regra de alerta à decisão da pessoa, em BPMN. Os números indicam os passos
deste capítulo.}
\label{fig:ia}
\end{figure}

\begin{remark}[info]
\textbf{Para TI.} O gatilho é a opção de pré-análise da regra de alerta (nível
\tech{agent\_preanalysis}). A entrega ao agente é feita pelo worker, com nova tentativa;
o resultado é buscado pelo Hanaro. Se o agente falha, só a pré-análise atrasa.
\end{remark}

\section{As peças do agente}

O agente é um serviço à parte do Hanaro, no servidor de IA. Ele não acessa o banco do Hanaro: só recebe os casos que o Hanaro entrega, responde e guarda a sua própria memória. A \autoref{fig:agente-componentes} mostra as seis peças do agente e como elas se ligam, no formato de um diagrama de componentes da UML. Cada caixa é uma peça; o círculo é o que a peça oferece, o meio círculo é o que ela precisa de outra, e as pequenas caixas na borda são as portas por onde o serviço conversa com o mundo de fora.

\begin{table}[H]
\caption{As peças do agente de IA.}\label{tab:agente-pecas}
\guiazebra
\small
\noindent\begin{tabularx}{\textwidth}{@{}L{3.2cm}X@{}}
    \toprule
    \guiacabtabela
    \thead{Peça} & \thead{O que faz} \\
    \midrule
    \textbf{API HTTP} & É a porta de entrada do serviço. Recebe os pedidos do Hanaro, protegidos por chave de API, para criar jobs, buscar casos e conversar. \\
    \textbf{Fila e worker} & Guarda os pedidos de pré-análise em uma fila e executa os pedidos em ordem. O mesmo pedido enviado duas vezes não vira duas pré-análises. \\
    \textbf{Grafo de pré-análise} & O roteiro da pré-análise: reunir os fatos, buscar casos parecidos, dar nota, escolher a rota, adaptar a proposta, conferir e finalizar. \\
    \textbf{Memória de casos} & Reúne as análises aprovadas e as encontra por significado e por palavras. A reordenação dos resultados é opcional, e cada análise aceita aumenta a memória. \\
    \textbf{Cliente de modelos} & Fala com o servidor de modelos: gera textos e vetores, e exige respostas em formato estruturado quando o roteiro precisa. \\
    \textbf{Chat de consulta} & Responde às perguntas da Consulta, com as fontes, usando só os casos da memória. \\
    \bottomrule
\end{tabularx}
\end{table}

\begin{remark}[info]
\textbf{Para TI.} A fila de jobs fica no PostgreSQL do próprio agente, o mesmo banco que guarda a memória de casos (com pgvector). O agente nunca chama o Hanaro: quem entrega os pedidos e busca os resultados é o worker do Hanaro.
\end{remark}

\begin{figure}[H]
\centering
\includegraphics[width=\dimexpr\textheight-3.2cm\relax,height=\textwidth,keepaspectratio,angle=90]{img/diagramas/componentes-agent-service}
\caption{Peças do agente de IA, em diagrama de componentes.}\label{fig:agente-componentes}
\end{figure}

\section{Rotas, selos e situações}

Cada pré-análise indica a \textbf{rota} que o agente seguiu, e a rota diz o quanto a
proposta pode ser aproveitada. A \autoref{tab:rotas} traz as três rotas; a
\autoref{tab:situacoes}, as situações em que uma pré-análise pode estar.

\begin{table}[H]
\caption{Rotas da pré-análise.}\label{tab:rotas}
\guiazebra
\small
\noindent\begin{tabularx}{\textwidth}{@{}L{3.2cm}X@{}}
    \toprule
    \guiacabtabela
    \thead{Rota} & \thead{O que significa} \\
    \midrule
    \textbf{Recorrência} & O item já apareceu em casos anteriores do Hanaro. A proposta
        parte do que foi feito antes, e o relatório pode se ligar ao plano já
        existente. \\
    \textbf{Caso parecido} & Há casos parecidos na memória. A proposta se apoia em até
        três deles, sem tratá-los como a mesma causa. \\
    \textbf{Sem precedente} & Não há nada parecido na memória. O assistente traz só os
        fatos, o checklist 4M e perguntas para quem analisa; qualquer causa mostrada é
        uma hipótese a confirmar. \\
\end{tabularx}
\end{table}

\begin{table}[H]
\caption{Situações de uma pré-análise.}\label{tab:situacoes}
\guiazebra
\small
\noindent\begin{tabularx}{\textwidth}{@{}L{3.9cm}X@{}}
    \toprule
    \guiacabtabela
    \thead{Situação} & \thead{O que significa} \\
    \midrule
    \textbf{Aguardando decisão} & Está pronta e ninguém decidiu ainda. É a que conta no
        menu e na lista. \\
    \textbf{Aplicada} & As decisões foram aplicadas: a análise, o relatório e as
        tarefas aceitas foram criados. \\
    \textbf{Aplicada em parte} & Parte das decisões foi aplicada; o restante não pôde
        ser. \\
    \textbf{Dispensada} & Alguém a descartou, com motivo. \\
    \textbf{Substituída} & O agente preparou uma versão mais nova, e esta ficou como
        histórico. \\
\end{tabularx}
\end{table}

O selo \textbf{Pré-análise pronta} aparece na tela \textbf{Alertas}, ao lado do alerta
que gerou a pré-análise, e leva direto a ela. Enquanto o agente trabalha, a lista mostra
\emph{Preparando…}; se houver versão nova, o aviso \emph{Nova versão em preparo}.

% ---------------------------------------------------------------------------
%  Os seis passos. Coordenadas dos marcadores: esquemas provisórios.
% ---------------------------------------------------------------------------
\passo[0.012/0.140, 0.182/0.090, 0.182/0.190, 0.182/0.360]{01-assistente-lista}{Abrir a lista do Assistente}{%
    A lista reúne as pré-análises preparadas para o seu trabalho. O contador ao lado do
    menu mostra quantas estão prontas e aguardando você, e é atualizado a cada minuto.
    \begin{elementos}
        \item \caminho{Assistente}, no menu lateral: abre a lista de pré-análises.
        \item \textbf{Seções}: \botao{Pré-análises}, \botao{Consulta},
              \botao{Memória} e \botao{Desempenho}. Aparecem só as que o seu perfil
              pode usar.
        \item \textbf{Abas e filtros}: \botao{Aguardando decisão},
              \botao{Decididas} e \botao{Todas}; filtros por gatilho, rota, linha e
              período.
        \item \textbf{Pré-análise}: rota, item, linha, valor e o gatilho que a
              disparou. Um clique abre a pré-análise (passo~2).
    \end{elementos}
}

\passo[0.182/0.120, 0.182/0.400, 0.182/0.660, 0.592/0.660]{02-preanalise-leitura}{Ler a pré-análise}{%
    Leia a pré-análise de cima para baixo antes de decidir. Ela mostra o que o
    assistente viu e de onde tirou cada sugestão.
    \begin{elementos}
        \item \textbf{Rota e \emph{O que aconteceu}}: a faixa da rota explica em que a
              proposta se apoia. Os fatos trazem item, linha, ocorrências, período,
              valor e a comparação com a mediana das janelas anteriores.
        \item \textbf{Proposta para a análise}: título, descrição, tipo de defeito,
              família da causa (4M), causa (hipótese) e contenção imediata.
        \item \textbf{Casos parecidos}: análises aprovadas de outros casos, com os
              motivos da semelhança. \botao{Não é parecido} faz o assistente dar menos
              peso ao caso nas próximas buscas.
        \item \textbf{Por que o assistente propõe isso} e \textbf{Fontes}: a
              justificativa de cada campo e a origem do dado (caso, relatório, regra ou
              dado do alerta).
    \end{elementos}
}

\begin{remark}
Parecido não é a mesma causa. Confira os casos citados e o histórico da linha antes de
aproveitar uma sugestão. Na rota \emph{Sem precedente}, a causa é só uma hipótese.
\end{remark}

\passo[0.182/0.120, 0.182/0.360, 0.182/0.600, 0.182/0.800]{03-preanalise-decisao}{Aceitar, editar ou rejeitar}{%
    Nenhuma sugestão é aplicada até você decidir sobre ela.
    \begin{elementos}
        \item \textbf{Cada campo}: \botao{Aceitar}, \botao{Editar} ou
              \botao{Rejeitar}. \botao{Desfazer} volta atrás e \botao{Aceitar tudo}
              aceita os campos de uma vez. Ao rejeitar, o campo \emph{Por quê?}
              (opcional) registra o motivo: \emph{Está errado}, \emph{Não se aplica},
              \emph{Repete outra coisa}, \emph{Já foi feito} ou \emph{Outro motivo}.
        \item \textbf{Tarefas propostas}: cada tarefa também é aceita, editada ou
              rejeitada. Tarefas só existem dentro de um relatório.
        \item \campo{Criar o relatório e finalizar a análise}: sem marcar esta opção, a
              análise fica como rascunho para você completar. Com ela, informe o título
              e o prazo do relatório. O responsável em branco fica com você.
        \item \botao{Aplicar decisões}: grava o que foi aceito ou editado.
              \botao{Dispensar} descarta a pré-análise, com um motivo.
    \end{elementos}
    Ao final, a pergunta \emph{A pré-análise ajudou?} (\emph{Sim}, \emph{Em parte},
    \emph{Não}) alimenta o Desempenho.
}

\begin{remark}[info]
Se a pré-análise cobre um grupo de ocorrências, marque \emph{Criar o relatório}: o
assistente cria uma análise para cada ocorrência que ainda não tem uma e um relatório
com todas, usando as mesmas decisões. Ocorrências que já tinham análise são mantidas.
\end{remark}

\begin{remark}
Se a tela avisar que \emph{há uma versão mais nova desta pré-análise} ou que ela
\emph{mudou enquanto você a revisava}, clique em \botao{Carregar a versão nova}. As
decisões feitas sobre a versão antiga não são aplicadas sozinhas.
\end{remark}

\passo[0.712/0.090, 0.182/0.780, 0.182/0.220, 0.182/0.490]{04-consulta}{Perguntar na Consulta}{%
    A Consulta responde a perguntas sobre como casos parecidos foram tratados. As
    respostas vêm só de análises já aprovadas e sempre mostram de onde saíram. Está
    disponível a Analistas e Gestores, quando a função está ligada.
    \begin{elementos}
        \item \botao{Histórico} e \botao{Nova conversa}: as conversas ficam salvas só
              para você. \botao{Apagar conversa} as remove de vez.
        \item \textbf{Pergunta}: escreva em texto livre, por exemplo \emph{Como
              resolvemos rebarba no chassi da linha G12?}, e clique em
              \botao{Perguntar}.
        \item \textbf{Resposta}: o texto do assistente, montado a partir dos casos
              encontrados.
        \item \textbf{Fontes}: cada análise citada, com linha, item e data.
              \botao{Abrir análise} leva até ela. \botao{Útil} e
              \botao{Não útil} avaliam a resposta.
    \end{elementos}
}

\begin{remark}[dica]
Sem resultado, tente outras palavras, o código do item ou a linha. A mensagem \emph{Nenhuma
análise parecida no histórico} indica que a memória não tem o que responder, e o
assistente não preenche essa lacuna.
\end{remark}

\passo[0.182/0.120, 0.182/0.290, 0.182/0.400, 0.182/0.800]{05-memoria}{Cuidar da memória e do desempenho}{%
    A \textbf{Memória} guarda as análises aprovadas que o assistente usa como exemplo.
    Os Gestores e os Administradores a mantêm.
    \begin{elementos}
        \item \textbf{Indicadores}: análises aprovadas, na memória, aguardando envio,
              retiradas, na fila de envio e envios com falha.
        \item \textbf{Busca e situação}: procure por título ou item e filtre por
              \emph{Na memória}, \emph{Aguardando envio} ou \emph{Retirada}.
        \item \textbf{Lista}: item, linha e data de cada análise.
              \botao{Retirar da memória} (com motivo opcional) tira a análise dos
              exemplos; \botao{Colocar de volta} a devolve.
        \item \botao{Reenviar tudo}: marca todas as análises para reenvio, em lotes e em
              segundo plano. Use depois de trocar o modelo de embedding ou se a
              memória se perder.
    \end{elementos}
}

\begin{remark}[info]
O \textbf{Desempenho} (\caminho{Assistente\seta Desempenho}) mostra, para um período, o
número de pré-análises decididas, a parcela das propostas aceitas como estavam, o tempo
mediano até a pré-análise ficar pronta, os campos mais rejeitados e seus motivos, a
avaliação de quem usou e o modelo em uso. Os números vêm do que as pessoas decidiram e
não dependem de o assistente estar ligado.
\end{remark}

\passo[0.212/0.140, 0.212/0.370, 0.212/0.520, 0.212/0.670, 0.562/0.840]{06-traducao}{Traduzir o relatório para o inglês}{%
    A tradução acontece ao emitir o documento de um relatório, em \botao{Baixar
    documento}.
    \begin{elementos}
        \item \campo{Idioma do documento}: escolha \emph{English}.
        \item \campo{Traduzir o texto para inglês}: opção com a etiqueta \emph{IA},
              que traduz o texto do documento antes do download. Só aparece quando a
              tradução está ligada.
        \item \textbf{Andamento}: \emph{Aguardando o serviço de tradução}, \emph{Na fila
              de tradução} (com o número de pedidos na sua frente), \emph{Carregando o
              modelo de tradução} e \emph{Traduzindo o texto, X de Y}.
        \item \textbf{Conferência}: se algum texto não passou na conferência, a tela
              lista os que ficarão em português. \botao{Emitir mesmo assim} segue com
              eles.
        \item \botao{Baixar PDF} (ou \botao{Baixar apresentação}): gera o documento.
    \end{elementos}
}

\begin{remark}[dica]
Se a tradução falhar, a tela oferece desligá-la e emitir em português, ou tentar de
novo. O documento em português continua disponível.
\end{remark}

\section{Como a tradução protege o seu relatório}

A tradução não entrega o relatório inteiro ao modelo. O Hanaro \textbf{mascara} antes
o que não pode mudar (números, códigos e marcadores), envia só o texto livre ao serviço
de tradução e, na volta, confere o resultado. Um campo só é gravado em inglês se passar
nas conferências de formato, marcadores e números. O campo que não passa fica em
português, e a tela avisa.

\begin{itemize}
    \item \textbf{Um pedido por vez.} Só um trabalho de tradução fala com o serviço de
          cada vez; os demais esperam na fila, do mais antigo para o mais novo.
    \item \textbf{Falha do serviço.} O Hanaro repete as falhas temporárias de rede. Se
          o serviço seguir indisponível, o trabalho termina como falho e nada
          incompleto é gravado.
    \item \textbf{Revisão humana preservada.} Ao trocar o modelo, as traduções
          automáticas antigas são refeitas, e as revisadas por pessoas permanecem.
\end{itemize}

\begin{remark}[info]
\textbf{Para TI.} A tradução usa o \tech{translation-service}, que fala com um servidor
de modelos de API compatível com a da OpenAI. Ativar, escolher o motor
(\tech{TRANSLATION\_ENGINE}) e trocar o modelo (incrementar
\tech{TRANSLATION\_ENGINE\_REVISION}).
\end{remark}

\section{Propostas de tarefas em um relatório}

Dentro de um relatório, o botão \botao{Propor tarefas} pede ao assistente sugestões de
tarefas, a partir do relatório e dos casos parecidos. Cada proposta mostra o motivo e as
fontes, e traz perguntas para quem analisa. Escolha as que fazem sentido e clique em
\botao{Criar as tarefas escolhidas}; \botao{Dispensar} descarta as propostas. Nada vira
tarefa até que você aceite.

\section{Boas práticas e limites}

\begin{itemize}
    \item \textbf{Leia antes de aceitar.} A proposta é um ponto de partida. Quem conhece
          a linha confirma a causa e a ação.
    \item \textbf{Rejeite com motivo.} O motivo informado alimenta o Desempenho e ajuda
          a equipe a ajustar o assistente.
    \item \textbf{Marque o que não é parecido.} \botao{Não é parecido} melhora as
          próximas buscas.
    \item \textbf{Cuide da memória.} Análises erradas ou pouco confiáveis podem ser
          retiradas dos exemplos; a memória só cresce com análises aprovadas.
    \item \textbf{Sem dado, sem resposta.} Se a memória não tem caso parecido, a
          Consulta diz isso e a pré-análise não inventa uma causa.
    \item \textbf{Modelo pequeno, conferência rígida.} O assistente usa modelos de
          linguagem de porte reduzido, e por isso tudo o que ele propõe passa por
          conferências: fonte inexistente, número ou data fora dos dados, família de
          causa ou tipo de defeito fora dos cadastros e texto acima dos limites do
          Hanaro são recusados.
\end{itemize}

\section{Mensagens e solução de problemas}

\begin{table}[H]
\caption{Mensagens do assistente e da tradução.}\label{tab:mensagens}
\guiazebra
\small
\noindent\begin{tabularx}{\textwidth}{@{}L{5.4cm}X@{}}
    \toprule
    \guiacabtabela
    \thead{Mensagem} & \thead{O que fazer} \\
    \midrule
    \emph{O assistente está desligado neste ambiente} & Siga o trabalho sem ele; as
        pré-análises já criadas continuam disponíveis. Peça à TI para ligá-lo. \\
    \emph{O assistente ainda está preparando esta pré-análise} & Aguarde e atualize a
        tela. \\
    \emph{O assistente não conseguiu preparar esta pré-análise} & Faça a análise normalmente
        na Base de Scrap; avise a TI se se repetir. \\
    \emph{Há uma versão mais nova desta pré-análise} & Clique em \botao{Carregar a
        versão nova} e decida de novo. \\
    \emph{Já existe uma análise desta ocorrência} & Abra a análise existente na Base de
        Scrap. \\
    \emph{O assistente não respondeu. Tente de novo em instantes} & A pergunta continua no
        campo da Consulta; envie de novo. \\
    \emph{O assistente está respondendo outras perguntas} & Aguarde alguns instantes e
        tente de novo. \\
    \emph{Nenhuma análise parecida no histórico} & Reescreva a pergunta com outras
        palavras, o código do item ou a linha. \\
    \emph{Não foi possível traduzir agora} & Desligue a tradução para emitir em
        português, ou tente de novo. \\
    \emph{N texto(s) ficarão em português} & Revise esses trechos ou use \botao{Emitir
        mesmo assim}. \\
\end{tabularx}
\end{table}

\fimdecapitulo


% ---------- FECHAMENTO ----------
\guiacontracapa

\end{document}
